\subsubsection{Penalty Terms}\label{sec:penalty-terms}

A \definition{Penalty Term} is an \textbf{additional contribution} to the objective function whose purpose is to \textbf{discourage solutions that violate a constraint}. The general form is:
\begin{equation}\label{eq:qubo-with-penalty}
    \argmin_{\mathbf{x}} \left( \mathbf{x}Q\mathbf{x}^{T} + \gamma \cdot \text{penalty}(\mathbf{x}) \right)
\end{equation}
Where the penalty function is designed such that:
\begin{equation}
    \text{penalty}(\mathbf{x}) = \begin{cases}
        0   & \text{if } \mathbf{x} \text{ satisfies the constraints} \\
        > 0 & \text{otherwise}
    \end{cases}
\end{equation}
And $\gamma$ controls the strength of the penalty.

\highspace
\begin{flushleft}
    \textcolor{Green3}{\faIcon{tools} \textbf{Common Constraints}}
\end{flushleft}
In the following, we will see some constraints and their corresponding penalty terms, in particular, we will see how to build the penalty term for the following constraints:
\begin{itemize}
    \begin{multicols}{2}
        \item $x_{1} + x_{2} \le 1$ (\autopageref{ex:penalty-terms-x1+x2-le-1})
        \item $x_{1} + x_{2} \ge 1$ (\autopageref{ex:penalty-terms-x1+x2-ge-1})
        \item $x_{1} + x_{2} = 1$ (\autopageref{ex:penalty-terms-x1+x2-eq-1})
        \item $x_{1} \le x_{2}$ (\autopageref{ex:penalty-terms-x1-le-x2})
        \item $x_{1} = x_{2}$ (\autopageref{ex:penalty-terms-x1-eq-x2})
        \item $x_{1} + x_{2} + x_{3} \le 1$ (\autopageref{ex:penalty-terms-x1+x2+x3-le-1})
    \end{multicols}
\end{itemize}

\highspace
\begin{examplebox}[: $x_{1} + x_{2} \le 1$]\phantomsection\label{ex:penalty-terms-x1+x2-le-1}
    Let's consider the \textbf{constraint}:
    \begin{equation*}
        x_{1} + x_{2} \le 1
    \end{equation*}
    Means that at most one variable can be 1. If we see at the truth table, we can see that the only way to violate the constraint is to have both $x_1$ and $x_2$ equal to 1:
    \begin{center}
        \begin{tabular}{@{} c c c c c @{}}
            \toprule
            $x_1$ & $x_2$ & $x_1 + x_2$ & Constraint Satisfied? & Penalty Desired \\
            \midrule
            0 & 0 & 0 & \textcolor{Green3}{\faIcon{check}} & 0 \\
            0 & 1 & 1 & \textcolor{Green3}{\faIcon{check}} & 0 \\
            1 & 0 & 1 & \textcolor{Green3}{\faIcon{check}} & 0 \\
            1 & 1 & 2 & \textcolor{Red2}{\faIcon{times}} & positive \\
            \bottomrule
        \end{tabular}
    \end{center}
    Thus, the only simple way to define a penalty term for this constraint is to multiply $x_1$ and $x_2$ together, since the product will be 1 only when both variables are 1, which is the only case that violates the constraint.

    The \textbf{penalty term} can be defined as:
    \begin{equation*}
        x_1 x_2 \qquad x_1 x_2 = \begin{cases}
            1 & \text{if } x_1 = x_2 = 1 \\
            0 & \text{otherwise}
        \end{cases}
    \end{equation*}
\end{examplebox}

\highspace
\begin{examplebox}[: $x_{1} + x_{2} \ge 1$]\phantomsection\label{ex:penalty-terms-x1+x2-ge-1}
    Let's consider the \textbf{constraint}:
    \begin{equation*}
        x_{1} + x_{2} \ge 1
    \end{equation*}
    Means that at least one variable must be 1. Let's write the thruth table for this constraint:
    \begin{center}
        \begin{tabular}{@{} c c c c c @{}}
            \toprule
            $x_1$ & $x_2$ & $x_1 + x_2$ & Constraint Satisfied? & Penalty Desired \\
            \midrule
            0 & 0 & 0 & \textcolor{Red2}{\faIcon{times}} & positive \\
            0 & 1 & 1 & \textcolor{Green3}{\faIcon{check}} & 0 \\
            1 & 0 & 1 & \textcolor{Green3}{\faIcon{check}} & 0 \\
            1 & 1 & 2 & \textcolor{Green3}{\faIcon{check}} & 0 \\
            \bottomrule
        \end{tabular}
    \end{center}
    The question becomes, \emph{can we build a polynomial that is 1 only when both variables are 0?} Yes, we can. Let's think it this way: we want a function that is \emph{true} only when both variables are 0. Logically, we can write this as:
    \begin{equation*}
        \text{penalty}(x_1, x_2) = \left(x_1 = 0\right) \land \left(x_2 = 0\right) = \lnot x_{1} \land \lnot x_{2}
    \end{equation*}
    Mathematically, a $\lnot$ can be represented as $1 - x_i$, since $x_i$ is a binary variable:
    \begin{center}
        \begin{tabular}{@{} c c @{}}
            \toprule
            $x$ & $1 - x$ \\
            \midrule
            0 & 1 \\
            1 & 0 \\
            \bottomrule
        \end{tabular}
    \end{center}
    Thus, we can substitute the $\lnot$ with $1 - x_i$:
    \begin{equation*}
    \text{penalty}(x_1, x_2) = (1 - x_1) \land (1 - x_2)
    \end{equation*}
    Finally, we can use the same trick as before to represent the $\land$ with a product:
    \begin{center}
        \begin{tabular}{@{} c c c @{}}
            \toprule
            $a$ & $b$ & $a \land b$ \\
            \midrule
            0 & 0 & 0 \\
            0 & 1 & 0 \\
            1 & 0 & 0 \\
            1 & 1 & 1 \\
            \bottomrule
        \end{tabular}
    \end{center}
    Thus, we can write the \textbf{penalty term} as:
    \begin{equation*}
        \text{penalty}(x_1, x_2) = (1 - x_1)(1 - x_2) = 1 - x_1 - x_2 + x_1 x_2
    \end{equation*}
\end{examplebox}

\highspace
\begin{examplebox}[: $x_{1} + x_{2} = 1$]\phantomsection\label{ex:penalty-terms-x1+x2-eq-1}
    Let's consider the \textbf{constraint}:
    \begin{equation*}
        x_{1} + x_{2} = 1
    \end{equation*}
    Means that exactly one variable must be 1. Let's write the thruth table for this constraint:
    \begin{center}
        \begin{tabular}{@{} c c c c c @{}}
            \toprule
            $x_1$ & $x_2$ & $x_1 + x_2$ & Constraint Satisfied? & Penalty Desired \\
            \midrule
            0 & 0 & 0 & \textcolor{Red2}{\faIcon{times}} & positive \\
            0 & 1 & 1 & \textcolor{Green3}{\faIcon{check}} & 0 \\
            1 & 0 & 1 & \textcolor{Green3}{\faIcon{check}} & 0 \\
            1 & 1 & 2 & \textcolor{Red2}{\faIcon{times}} & positive \\
            \bottomrule
        \end{tabular}
    \end{center}
    Differently from the previous cases, here we have two cases that violate the constraint: when both variables are 0 and when both variables are 1. Thus, we need to find a function that is 1 in both cases. Again, we can use the logical representation to find the function:
    \begin{equation*}
        \left[
            \left(x_1 = 0\right) \land \left(x_2 = 0\right)
        \right]
        \lor
        \left[
            \left(x_1 = 1\right) \land \left(x_2 = 1\right)
        \right]
    \end{equation*}
    We already know how to represent the $\land$ and the $\lnot$, but we also need to represent the $\lor$. The $\lor$ truth table is the following:
    \begin{center}
        \begin{tabular}{@{} c c c @{}}
            \toprule
            $a$ & $b$ & $a \lor b$ \\
            \midrule
            0 & 0 & 0 \\
            0 & 1 & 1 \\
            1 & 0 & 1 \\
            1 & 1 & 1 \\
            \bottomrule
        \end{tabular}
    \end{center}
    At first, it might seem that we cannot represent the $\lor$ with a polynomial, since the only case that is 0 is when both variables are 0. However, we can use the \textbf{dual form of De Morgan's law for disjunction}:
    \begin{equation*}
        a \lor b = \lnot \left(\lnot a \land \lnot b\right)
    \end{equation*}
    And the truth table for this expression is the following:
    \begin{center}
        \begin{tabular}{@{} c c c c c c @{}}
            \toprule
            $a$ & $b$ & $\lnot a$ & $\lnot b$ & $\lnot a \land \lnot b$ & $\lnot \left(\lnot a \land \lnot b\right)$ \\
            \midrule
            0 & 0 & 1 & 1 & 1 & 0 \\
            0 & 1 & 1 & 0 & 0 & 1 \\
            1 & 0 & 0 & 1 & 0 & 1 \\
            1 & 1 & 0 & 0 & 0 & 1 \\
            \bottomrule
        \end{tabular}
    \end{center}
    Now we can substitute the $\lnot$ with $1 - x_i$ and the $\land$ with a product:
    \begin{equation*}
        \text{penalty}(x_1, x_2) = \left[
            \left(1 - x_1\right)\left(1 - x_2\right)
        \right]
        \lor
        \left[
            x_1 x_2
        \right]
    \end{equation*}
    Finally, we can substitute the $\lor$ with the dual form of De Morgan's law:
    \begin{equation*}
        \text{penalty}(x_1, x_2) = 1 - \left[
            1 - \left(1 - x_1\right)\left(1 - x_2\right)
        \right]
        \cdot
        \left[
            1 - x_1 x_2
        \right]
    \end{equation*}
    Which can be simplified to:
    \begin{align*}
        \text{penalty}(x_1, x_2) & = 1 - \left[
            1 - \left(1 - x_1\right)\left(1 - x_2\right)
        \right]
        \cdot
        \left[
            1 - x_1 x_2
        \right] \\
        & = 1 - \left[
            1 - \left(1 -x_2 -x_1 +x_1 x_2\right)
        \right]
        \cdot
        \left[
            1 - x_1 x_2
        \right] \\
        & = 1 - \left[
            x_2 + x_1 - x_1 x_2
        \right]
        \cdot
        \left[
            1 - x_1 x_2
        \right] \\
        & = 1 - \left(
            x_{2} -x_{1}x_{2}^{2}  +x_{1} -x_{1}^{2}x_{2}  -x_{1}x_{2} +x_{1}^{2}x_{2}^{2}
        \right) \\
        & = 1 -x_{2} +x_{1}x_{2}^{2}  -x_{1} +x_{1}^{2}x_{2}  +x_{1}x_{2} -x_{1}^{2}x_{2}^{2}
        \\
    \end{align*}
    Now, it might seem so complicated, but we can simplify it even more. Since $x_i$ is a binary variable, we can substitute $x_i^2$ with $x_i$:
    \begin{align*}
        \text{penalty}(x_1, x_2) & = 1 -x_{2} +x_{1}x_{2}^{2}  -x_{1} +x_{1}^{2}x_{2}  +x_{1}x_{2} -x_{1}^{2}x_{2}^{2} \\
        & = 1 -x_{2} +x_{1}x_{2}  -x_{1} +x_{1}x_{2}  +x_{1}x_{2} -x_{1}x_{2} \\
        & = 1 -x_{1} -x_{2} +2x_{1}x_{2}
    \end{align*}
    We got it! Let's check the truth table to be sure:
    \begin{itemize}
        \item $\text{penalty}(0, 0) = 1 - 0 - 0 + 2(0)(0) = 1$ (constraint not satisfied)
        \item $\text{penalty}(0, 1) = 1 - 0 - 1 + 2(0)(1) = 0$ (constraint satisfied)
        \item $\text{penalty}(1, 0) = 1 - 1 - 0 + 2(1)(0) = 0$ (constraint satisfied)
        \item $\text{penalty}(1, 1) = 1 - 1 - 1 + 2(1)(1) = 1$ (constraint not satisfied)
    \end{itemize}
    Everything checks out! So the \textbf{penalty term} can be defined as:
    \begin{equation*}
        \text{penalty}(x_1, x_2) = 1 -x_{1} -x_{2} +2x_{1}x_{2}
    \end{equation*}
\end{examplebox}

\highspace
\begin{examplebox}[: $x_{1} \le x_{2}$]\phantomsection\label{ex:penalty-terms-x1-le-x2}
    Let's consider the \textbf{constraint}:
    \begin{equation*}
        x_{1} \le x_{2}
    \end{equation*}
    The constraint means that $x_1$ can be 1 only if $x_2$ is also 1, otherwise, if $x_1$ is 0, $x_2$ can be either 0 or 1. Let's write the truth table for this constraint:
    \begin{center}
        \begin{tabular}{@{} c c c c @{}}
            \toprule
            $x_1$ & $x_2$ & Constraint Satisfied? & Penalty Desired \\
            \midrule
            0 & 0 & \textcolor{Green3}{\faIcon{check}} & 0 \\
            0 & 1 & \textcolor{Green3}{\faIcon{check}} & 0 \\
            1 & 0 & \textcolor{Red2}{\faIcon{times}} & positive \\
            1 & 1 & \textcolor{Green3}{\faIcon{check}} & 0 \\
            \bottomrule
        \end{tabular}
    \end{center}
    The only case that violates the constraint is when $x_1$ is 1 and $x_2$ is 0. Thus, we need to find a function that is 1 only in this case. We can use the logical representation to find the function:
    \begin{equation*}
        \text{penalty}(x_1, x_2) = \left(x_1 = 1\right) \land \left(x_2 = 0\right) = x_1 \land \lnot x_2
    \end{equation*}
    Now we can substitute the $\lnot$ with $1 - x_2$:
    \begin{equation*}
        \text{penalty}(x_1, x_2) = x_1 \land (1 - x_2)
    \end{equation*}
    Finally, we can substitute the $\land$ with a product:
    \begin{equation*}
        \text{penalty}(x_1, x_2) = x_1 (1 - x_2) = x_1 - x_1 x_2
    \end{equation*}
    Let's check the truth table to be sure:
    \begin{itemize}
        \item $\text{penalty}(0, 0) = 0 - 0 = 0$ (constraint satisfied)
        \item $\text{penalty}(0, 1) = 0 - 0 = 0$ (constraint satisfied)
        \item $\text{penalty}(1, 0) = 1 - 0 = 1$ (constraint not satisfied)
        \item $\text{penalty}(1, 1) = 1 - 1 = 0$ (constraint satisfied)
    \end{itemize}
    Everything checks out! So the \textbf{penalty term} can be defined as:
    \begin{equation*}
        \text{penalty}(x_1, x_2) = x_1 - x_1 x_2
    \end{equation*}
\end{examplebox}

\highspace
\begin{examplebox}[: $x_{1} = x_{2}$]\phantomsection\label{ex:penalty-terms-x1-eq-x2}
    Let's consider the \textbf{constraint}:
    \begin{equation*}
        x_{1} = x_{2}
    \end{equation*}
    The constraint means that both variables must be the same. Let's write the truth table for this constraint:
    \begin{center}
        \begin{tabular}{@{} c c c c @{}}
            \toprule
            $x_1$ & $x_2$ & Constraint Satisfied? & Penalty Desired \\
            \midrule
            0 & 0 & \textcolor{Green3}{\faIcon{check}} & 0 \\
            0 & 1 & \textcolor{Red2}{\faIcon{times}} & positive \\
            1 & 0 & \textcolor{Red2}{\faIcon{times}} & positive \\
            1 & 1 & \textcolor{Green3}{\faIcon{check}} & 0 \\
            \bottomrule
        \end{tabular}
    \end{center}
    The only cases that violate the constraint are when $x_1$ is 0 and $x_2$ is 1, or when $x_1$ is 1 and $x_2$ is 0. Thus, we need to find a function that is 1 in both cases. We can use the logical representation to find the function:
    \begin{equation*}
        \text{penalty}(x_1, x_2) = \left[
            \left(x_1 = 0\right) \land \left(x_2 = 1\right)
        \right]
        \lor
        \left[
            \left(x_1 = 1\right) \land \left(x_2 = 0\right)
        \right]
    \end{equation*}
    Now we can substitute the $\lnot$ with $1 - x_i$ and the $\land$ with a product:
    \begin{align*}
        &= \left[
            (1 - x_1) x_2
        \right]
        \lor
        \left[
            x_1 (1 - x_2)
        \right] \\
        & = 1 - \left[
            1 - (1 - x_1) x_2
        \right]
        \cdot
        \left[
            1 - x_1 (1 - x_2)
        \right] \\
        & = 1 - \left[
            1 - x_2 + x_1 x_2
        \right]
        \cdot
        \left[
            1 - x_1 + x_1 x_2
        \right] \\
        & = 1 - \left(
            1 - x_1 + x_1 x_2 %
            - x_2 + x_1 x_2 - x_1 x_2^2 %
            + x_1 x_2 - x_1^2 x_2 + x_1^2 x_2^2
        \right) \\
        & \downarrow \text{ since } x_i^2 = x_i \\
        & = 1 - \left(
            1 - x_1 + x_1 x_2 %
            - x_2 + x_1 x_2 - x_1 x_2 %
            + x_1 x_2 - x_1 x_2 + x_1 x_2
        \right) \\
        & \downarrow \text{ simplify} \\
        & = 1 - \left(
            1 - x_1 + 2 x_1 x_2 - x_2
        \right) \\
        & = x_1 + x_2 - 2 x_1 x_2
    \end{align*}
    Let's check the truth table to be sure:
    \begin{itemize}
        \item $\text{penalty}(0, 0) = 0 + 0 - 0 = 0$ (constraint satisfied)
        \item $\text{penalty}(0, 1) = 0 + 1 - 0 = 1$ (constraint not satisfied)
        \item $\text{penalty}(1, 0) = 1 + 0 - 0 = 1$ (constraint not satisfied)
        \item $\text{penalty}(1, 1) = 1 + 1 - 2 = 0$ (constraint satisfied)
    \end{itemize}
    Everything checks out! So the \textbf{penalty term} can be defined as:
    \begin{equation*}
        \text{penalty}(x_1, x_2) = x_1 + x_2 - 2 x_1 x_2
    \end{equation*}
\end{examplebox}

\highspace
\begin{examplebox}[: $x_{1} + x_{2} + x_{3} \le 1$]\phantomsection\label{ex:penalty-terms-x1+x2+x3-le-1}
    Let's consider the \textbf{constraint}:
    \begin{equation*}
        x_{1} + x_{2} + x_{3} \le 1
    \end{equation*}
    The constraint means that at most one variable can be 1. Let's write the truth table for this constraint:
    \begin{center}
        \begin{tabular}{@{} c c c c c @{}}
            \toprule
            $x_1$ & $x_2$ & $x_3$ & $x_1 + x_2 + x_3$ & Constraint Satisfied? \\
            \midrule
            0 & 0 & 0 & 0 & \textcolor{Green3}{\faIcon{check}} \\
            0 & 0 & 1 & 1 & \textcolor{Green3}{\faIcon{check}} \\
            0 & 1 & 0 & 1 & \textcolor{Green3}{\faIcon{check}} \\
            0 & 1 & 1 & 2 & \textcolor{Red2}{\faIcon{times}} \\
            1 & 0 & 0 & 1 & \textcolor{Green3}{\faIcon{check}} \\
            1 & 0 & 1 & 2 & \textcolor{Red2}{\faIcon{times}} \\
            1 & 1 & 0 & 2 & \textcolor{Red2}{\faIcon{times}} \\
            1 & 1 & 1 & 3 & \textcolor{Red2}{\faIcon{times}} \\
            \bottomrule
        \end{tabular}
    \end{center}
    The constraint is satisfied when at most one variable is set to true (i.e., equals to one). Here, instead of trying to find a logical representation of the penalty term, we can think more abstractly. Invalid states are all cases where \textbf{at least two variables are 1}:
    \begin{equation*}
        110, \, 101, \, 011, \, 111
    \end{equation*}
    From this, we can see that $110$ is simply:
    \begin{equation*}
        x_1 x_2 (1 - x_3)
    \end{equation*}
    But we could include $111$ as well, since it also violates the constraint. Thus, we can write something like $11-$ where the $-$ means that the variable can be either 0 or 1, don't care about it. This can be represented as:
    \begin{equation*}
        \left(x_1 = 1\right) \land \left(x_2 = 1\right) \land - = x_1 x_2
    \end{equation*}
    Now, the remaining invalid states are $101$ and $011$, which can be represented as:
    \begin{equation*}
        101 \to \left(x_1 = 1 \land x_2 = 0 \land x_3 = 1\right) \qquad 011 \to \left(x_1 = 0 \land x_2 = 1 \land x_3 = 1\right)
    \end{equation*}
    But again, we don't care if $x_2$ is 0 or 1 in $101$, and similarly, we don't care if $x_1$ is 0 or 1 in $011$. Thus, we can simply write:
    \begin{equation*}
        101 \to \left(x_1 = 1 \land x_3 = 1\right) \qquad 011 \to \left(x_2 = 1 \land x_3 = 1\right)
    \end{equation*}
    Which can be represented as:
    \begin{equation*}
        101 \to x_1 x_3 \qquad 011 \to x_2 x_3
    \end{equation*}
    Finally, we can combine all the invalid states together to get the penalty term:
    \begin{equation*}
        \text{penalty}(x_1, x_2, x_3) = x_1 x_2 + x_1 x_3 + x_2 x_3
    \end{equation*}
    Let's check the truth table to be sure:
    \begin{itemize}
        \item $\text{penalty}(0, 0, 0) = 0 + 0 + 0 = 0$ (constraint satisfied)
        \item $\text{penalty}(0, 0, 1) = 0 + 0 + 0 = 0$ (constraint satisfied)
        \item $\text{penalty}(0, 1, 0) = 0 + 0 + 0 = 0$ (constraint satisfied)
        \item $\text{penalty}(0, 1, 1) = 0 + 0 + 1 = 1$ (constraint not satisfied)
        \item $\text{penalty}(1, 0, 0) = 0 + 0 + 0 = 0$ (constraint satisfied)
        \item $\text{penalty}(1, 0, 1) = 0 + 1 + 0 = 1$ (constraint not satisfied)
        \item $\text{penalty}(1, 1, 0) = 1 + 0 + 0 = 1$ (constraint not satisfied)
        \item $\text{penalty}(1, 1, 1) = 1 + 1 + 1 = 3$ (constraint not satisfied)
    \end{itemize}
    Everything checks out! So the \textbf{penalty term} can be defined as:
    \begin{equation*}
        \text{penalty}(x_1, x_2, x_3) = x_1 x_2 + x_1 x_3 + x_2 x_3
    \end{equation*}
\end{examplebox}

\highspace
In general, a constraint is converted into a QUBO penalty by building a polynomial that is:
\begin{equation*}
    \text{penalty}(\mathbf{x}) = 0
\end{equation*}
For all \textbf{feasible assignments}, and:
\begin{equation*}
    \text{penalty}(\mathbf{x}) > 0
\end{equation*}
For all \textbf{infeasible assignments}. Therefore, \textbf{constraints} are not imposed explicitly: they are transformed into \textbf{energy costs}. The optimizer is still free to choose any binary assignment, but \hl{infeasible assignments become energetically unfavorable}.

\highspace
A useful way to construct penalty terms is:
\begin{enumerate}
    \item Identify the \textbf{assignments that violate} the constraint;
    \item \textbf{Build a detector polynomial} for each forbidden assignment or forbidden pattern;
    \item \textbf{Sum all detector polynomials} to get the overall penalty function;
    \item \emph{Optionally} multiply the result by a penalty coefficient $\gamma > 0$.
\end{enumerate}
For binary variables, the basic translations are:
\begin{gather*}
    x_i = 1 \;\to\; x_i,
    \qquad
    x_i = 0 \;\to\; 1-x_i,
    \\
    \text{AND} \;\to\; a \land b \;\to\; ab,
    \qquad
    \text{OR} \;\to\; \lnot(\lnot x_1 \land \lnot x_2) \; \to\; 1 - (1-x_1)(1-x_2)
\end{gather*}
Thus, the \textbf{penalty term acts as an energy detector} for constraint violations: feasible solutions have zero extra energy, while infeasible solutions receive a positive energy penalty.